Work first on the compact quotient with boundary $\widetilde\Omega=\Omega\times(\mathbb R/\mathbb Z)^m$, extending the data constantly in the new periodic directions. In the four-slot basis $(1,\mathbf i,\mathbf j,\mathbf k)$ with indices $0,1,2,3$, the original periods $\Lambda$ are in slot $0$ and the new periods $\mathbb Z^m$ are in slot $3$. Regrouping the real coordinates identifies $\widetilde\Omega$ with the torus $\mathbb R^{2m}/(\Lambda\times\mathbb Z^m)$ in slots $0,3$ times $V$ in slots $1,2$. Until descent, $G,H$ are the full four-slot gradient and Hessian. Write $\mathcal H(D^2\psi)=H(\psi)$, and write $\widehat\xi$ for a real covector $\xi$ encoded as a quaternionic column by its four slots. Lemma~\ref{ms:lemma-3-1} gives $\mathcal H(\xi\xi^{\mathsf t})=\widehat\xi\widehat\xi^*$, $\operatorname{tr}H(\psi)=\Delta\psi$, and covariance under quaternionic unitary coordinate transformations.

\emph{Comparison and continuation.} Join the known solution $u=r$ to the desired density by
\[
 g_t=(1-t)\det H(B_0+r)+tg,\qquad 0\le t\le1.
\]
At a solution $\psi=B_0+u$, put $Q=H(\psi)$ and $S=\operatorname{tr}Q^{-1}$. By Lemma~\ref{ms:lemma-3-2}, the derivative of the log determinant and its real principal symbol are
\[
 Lw=Q^{-1}:H(w),\qquad \sigma_L(\xi)=\widehat\xi^*Q^{-1}\widehat\xi .
\]
Thus $L$ is uniformly elliptic at any fixed smooth solution with $Q>0$ on the compact closure. If $Q_i=H(\psi_i)>0$ for two potentials, interpolate the matrices to obtain
\[
 \log\det Q_1-\log\det Q_2=
 \left(\int_0^1((1-\theta)Q_2+\theta Q_1)^{-1}\,\mathrm d\theta\right):H(\psi_1-\psi_2).
\]
The averaged real operator is uniformly elliptic when the Hessians are positive up to the boundary. The maximum principle proves the comparison assertion in four slots and implies uniqueness. The same argument holds for a bounded subdomain or a compact quotient with boundary.

For each $k\ge0$ and $0<\alpha<1$, the linearized Dirichlet operator $L:C_0^{k+2,\alpha}\to C^{k,\alpha}$ is an isomorphism, where the subscript means zero boundary trace. The local linear Schauder estimates and Dirichlet theory of \cite{r-31}, Theorems 6.14 and 6.19, apply in a finite cover of the flat quotient. In the method of continuity from the Laplacian to $L$, the defining function $r$ is a common strict barrier since $H(r)\ge cI$. The maximum principle gives the uniform $C^0$ estimate for this linear interpolation, the Schauder estimates give the corresponding higher bounds, and the initial Dirichlet Laplacian is invertible. The implicit function theorem applied to the nonlinear equation now proves openness of the set of solvable $t$.

With this product and the four-slot Hessian, Lemma~\ref{ms:dirichlet-c2} provides the two barriers and the uniform full $C^2$ estimate for each smooth solution on this path.

\emph{Regularity, descent and dependence on parameters.} The $C^2$ bound and $\det Q=g_t\ge c_0>0$ give $cI\le Q\le CI$. Thus the real Hessians lie in the compact convex set
\[
 \mathcal K=\{M:|M|\le C_1,\ cI\le\mathcal H(M)\le CI\}.
\]
On $\mathcal H(M)>0$ the smooth function $F(M)=\log\det\mathcal H(M)$ is concave by Lemma~\ref{ms:lemma-3-2}. For $M\in\mathcal K$ one has
\[
 DF(M)[\xi\xi^{\mathsf t}]=
 \widehat\xi^*\mathcal H(M)^{-1}\widehat\xi
 \in[C^{-1}|\xi|^2,c^{-1}|\xi|^2].
\]
These are the uniform real ellipticity and concavity hypotheses on $\mathcal K$ in the interior and boundary estimates of \cite{r-31}, Theorems 17.14 and $17.26'$, hypotheses (i)$'$ and (ii)$'$. To define the operator on the full real matrix space, multiply $F$ by a smooth cutoff supported where $\mathcal H(M)>0$ and equal to one on a neighborhood of $\mathcal K$. The equation for the estimate is $F(D^2\psi)-\log g_t=0$ with smooth boundary value $B_0$. Applying the local estimates in a finite covering gives a uniform $C^{2,\alpha}$ bound for some $\alpha>0$. At a limit of solvable parameters, extract a subsequence converging in $C^{2,\beta}$ for $0<\beta<\alpha$. Positivity and the equation pass to the limit. Schauder bootstrapping yields smoothness up to the boundary: after flattening, differentiate tangentially with the specified Dirichlet data and recover the pure normal second derivative from the equation by uniform ellipticity. This proves closedness, so the path reaches $g$.

Translations in the new periodic coordinates preserve the data and domain. Uniqueness on $\widetilde\Omega$ makes the solution invariant under all such translations, so it descends to $\Omega$. Any physical solution lifts, which gives uniqueness, and the same lift gives comparison on physical subdomains. For a smooth family, pull the lifted problem back to a fixed domain. Its coefficients and boundary values depend smoothly on the parameters, and the linearization is conjugate to the Dirichlet isomorphism. The implicit function theorem in each $C_0^{k+2,\alpha}$ proves smooth dependence up to the boundary. The strict determinant inequalities choosing $b,b'$ persist locally for the family.
